\makeabstract
{
Signals: Searching for Long-Range Spin-Dependent Interactions
}
{
Yana Boneva (1), Dr. Nathan Clayburn (1,2), Andrew Glassford (1), Ashley Kim (1), Thomas Uelmen (1), Professor Larry Hunter (1)
}
{
\textsuperscript{1}Department of Physics and Astronomy, Amherst College, Amherst, MA, USA\\
\textsuperscript{2}Department of Physics, Keene State College, Keene, NH, USA
}
{
Theorized extensions of the Standard Model predict an additional, spin-dependent fifth fundamental force. Our experiment uses the Earth as a source of spin-polarized electrons, and an 199Hg-133Cs co-magnetometer as a detector to look for long-range candidate particles for that force. The electron spin in Cs has been shown to not exhibit any exotic behavior, so we use the Cs signal as a baseline for the local magnetic fields and their associated effects. This leaves the nuclear spin of 199Hg and its associated valence neutron sensitive to exotic spin-dependent forces. The units of our measurement are μHz of Hg precession frequency (Hg μHz).
}
{
Our set-up uses five magnetic field coils to create a homogeneous magnetic field within the apparatus, and an additional, sixth coil to cancel out Earth’s field. The Hg and Cs atoms are first pumped up with circularly polarized light tuned to their absorption frequency and intensity modulated at their Larmor precession frequency. They are then probed with linearly polarized light tuned away from their absorption frequency. A rotating table in our apparatus aligns the precession field of the atoms with north, south, west and east to probe the different alignments to dot or cross product-dependent candidate interactions. The North-South and East-West signals are then subtracted to get our final result. Our data is collected into octets, EW or NS, consisting of 8 distinct table positions ordered as ABBABAAB to allow us to subtract off first and second order drifts. The octets are then divided into sagas based on their time sequence and magnetic field orientation. The direction of the precession fields is flipped for half of the data to measure and cancel out systematic effects related to table position.
}
{
A null result in NS is 17.112 Hg μHz and 0.000 Hg μHz in EW. We found results of -0.53 ± 0.056 stat  ± 0.04 syst μHz in EW and 17.72 ± 0.055 stat  ± 0.04 syst μHz in NS.
}
{
We achieved our aim of 0.07 Hg μHz of systematic and statistical uncertainty. This experiment will place bounds on these candidate interactions one order of magnitude better than our laboratory’s previous experiment.
}

\newpage
